% byu-coi-plan.tex
%
% THE INSTRUMENT. This is the document that gets attached to BYU's Financial
% Disclosure Form, approved, and signed. It states obligations only.
%
% The reasons, the approval sequence, the signing tiers, and the company
% formation-document analysis live in byu-coi-briefing.tex, which is NOT
% submitted. If you are about to add a sentence here that explains WHY a
% clause is justified, it belongs in the briefing or in a footnote.
%
% Hard limits, both checked at build time:
%   * six pages maximum
%   * no em or en dashes (see the house rule in vcl-doc.sty)
%
% Compile with pdflatex, THREE times, then copy the PDF to the repo root
% (latex/.gitignore excludes *.pdf, so the tracked artifact is ../*.pdf).
%
\documentclass[11pt]{article}
\usepackage{vcl-doc}

\hypersetup{
  pdftitle={Conflict of Interest Management Plan},
  pdfsubject={Conflict-of-interest management plan for a faculty co-founder},
  pdfkeywords={BYU, conflict of interest, management plan, faculty founder}
}
\vclrunninghead{Conflict of Interest Management Plan, Vertical Cloud Lab}

\begin{document}

% ================================================================== title ==
\thispagestyle{plain}

{\color{byunavy}\hrule height 2.4pt}
\vspace{12pt}
{\LARGE\bfseries\color{byunavy} Conflict of Interest Management Plan}\\[6pt]
{\normalsize\bfseries Brigham Young University}\\[2pt]
{\normalsize Department of Mechanical Engineering, Ira A.\ Fulton College of
Engineering}
\vspace{8pt}

{\color{byunavy}\hrule height 0.8pt}
\vspace{9pt}

{\setlength{\extrarowheight}{2pt}\small
\begin{tabular}{@{} >{\leavevmode\color{byunavy}\scshape}p{1.3in} >{\RaggedRight\arraybackslash}p{5.15in} @{}}
Faculty member & [Sterling G.\ Baird], Assistant Professor, Mechanical
Engineering (the ``Faculty Member'') \\
Company & [NewCo, Inc.], a [Delaware corporation] commercializing automated
(``cloud lab'') materials-experimentation technology (the ``Company'') \\
Governing policies &
\href{https://policy.byu.edu/view/conflict-of-time-or-interest-policy}{Conflict
of Time or Interest};
\href{https://policy.byu.edu/view/financial-conflict-of-interest-in-research-policy}{Financial
Conflict of Interest in Research};
\href{https://policy.byu.edu/view/intellectual-property-policy}{Intellectual
Property};
\href{https://spo.byu.edu/policies}{Financial Disclosure Form and
Instructions};
\href{https://finserve.byu.edu/00000193-6566-d995-a9db-eff6a47a0001/sponsored-programs-handbook}{Sponsored
Projects Handbook};
\href{https://spo.byu.edu/00000190-3687-d777-a591-37af392d0001/rechargecenterpolicyprocedures-pdf}{Recharge
Center Policy} \\
Effective date & [date] \\
Review & annual, with the Faculty Profile outside-activity disclosure \\
Companion & a separate briefing carries the approval sequence and the
reasoning behind these conditions; it is not part of this plan \\
\end{tabular}}



% ------------------------------------------------------------------- band A
\planband{A. Role and principal duties of the conflicted investigator}{}

\psection{The relationship and its limits}

\clause{1.1} The Faculty Member co-founded the Company, holds \textbf{more
than ten percent} of it (currently \ph\%, counting stock, options, rights to
acquire stock, and immediate family interests), and remains on a full-time
faculty appointment, not on leave.

\clause{1.2} \textbf{Advisory roles only.} The only Company roles are
stockholder and \textbf{[Chief Scientific Advisor]}: advisory, with no officer
title, no employment relationship, no board vote, and no authority over
Company staff.\footnote{\href{https://finserve.byu.edu/00000193-6566-d995-a9db-eff6a47a0001/sponsored-programs-handbook}{Sponsored
Projects Handbook}: ``Faculty members not on leave are not to be line officers
in outside companies.'' Also
\href{https://policy.byu.edu/view/conflict-of-time-or-interest-policy}{Conflict
of Time or Interest Policy}.}

\clause{1.3} \textbf{No pay from the Company}: no salary, consulting fees,
bonuses, honoraria, or other compensation. Stock is the only financial
interest.\footnote{\href{https://policy.byu.edu/view/intellectual-property-policy}{Intellectual
Property Policy}: a faculty member may consult for a company or take part in
research it sponsors, not both. Exceptions need chair, dean, and the written
approval of the Associate Academic Vice President for Research (``the AAVP'').}

\clause{1.4} \textbf{Travel is reimbursed at cost, and reported.} The Company
may reimburse the Faculty Member's documented airfare, lodging, and meals for
a trip whose purpose is Company business: actual cost, the Faculty Member
alone, no per diem above cost, nothing else of more than nominal value. Those days count against the four-day allowance (2.1); no trip
is paid twice, and university business, including conference travel to present
BYU research, stays on BYU funds; and each reimbursed trip is reported on the
\href{https://spo.byu.edu/policies}{annual financial disclosure} by sponsor,
purpose, dates, and destination, which is what that form asks for.

\clause{1.5} \textbf{The inventor's royalty share is forfeited.} At ten
percent or more the Faculty Member is a ``principal equity owner'' and gives
up the inventor's 45\% share of BYU's licensing income from the Company, with
the matched research-account option. Policy provides no
exception.\footnote{\href{https://policy.byu.edu/view/intellectual-property-policy}{Intellectual
Property Policy}: ``a principal equity owner is one who controls ten percent
or more equity in the company,'' and such a developer ``will forfeit his or
her distribution.''}

\clause{1.6} \textbf{No voting board seat}, and no observer seat unless the
written authorization of the outside role describes and authorizes one before
it is taken.\footnote{A voting seat is allowed only where the company ``has
not raised more than \$3,000,000,'' and then only ``on rare occasions'' with
the dean, the AAVP, and the Technology Transfer Office approving in writing
(\href{https://policy.byu.edu/view/intellectual-property-policy}{Intellectual
Property Policy}). The Company is past that ceiling, and policy does not
address observer seats at all.}

\clause{1.7} \textbf{Neither role speaks for the other.} The Faculty Member
does not represent BYU in any Company matter or the Company in any BYU matter,
and commits neither to anything. The Company does not use BYU's name, logo, or
facilities in its advertising or investor
materials.\footnote{\href{https://policy.byu.edu/view/intellectual-property-policy}{Intellectual
Property Policy}: ``The name of the university shall not be used\ldots\ in any
advertising.''}

% ------------------------------------------------------------------- band B
\planband{B. Conditions of the plan}{}

\psection{Time, teaching, and availability to students}

\clause{2.1} All Company activity fits inside BYU's outside-activity
allowance: \textbf{at most four days a month} across every outside activity
combined, with the chair's prior written
approval.\footnote{\href{https://policy.byu.edu/view/conflict-of-time-or-interest-policy}{Conflict
of Time or Interest Policy}: ``up to four days per month.''} The Faculty
Member logs days spent and gives the log to the chair with the annual
disclosure.

\clause{2.2} \textbf{Teaching and students come first.} Classes, class
preparation, and office hours are held as scheduled, and student and
research-group meetings go on the calendar before Company commitments. Company
time taken during a workday is made up on evenings or weekends that same
week.\footnote{\href{https://policy.byu.edu/view/conflict-of-time-or-interest-policy}{Conflict
of Time or Interest Policy}, compensating-time expectation.}

\clause{2.3} The chair reviews the log at annual review against the policy's
standard (availability to students; teaching, scholarship, and citizenship
performance) and may require Company time to be reduced.

\psection{Contracts: who negotiates and who signs}

\clause{3.1} The Faculty Member takes \textbf{no part in negotiating,
approving, or influencing any agreement between BYU and the Company, on either
side of the table}: the license, every research agreement, every equipment
placement, and any gift. The Company is represented by a non-founder officer
([name/title]); BYU by the Research Administration Office, the Technology
Transfer Office, or the Office of General
Counsel.\footnote{\href{https://policy.byu.edu/view/conflict-of-time-or-interest-policy}{Conflict
of Time or Interest Policy} lists as unmanageable ``negotiating, influencing,
or attempting to influence the negotiations of contracts or agreements between
the university and a third party for the benefit of the employee or a Related
Party.''}

\clause{3.2} The Faculty Member does not advocate for any such deal or its
terms with the chair, dean, Research Administration Office, Technology
Transfer Office, or the AAVP, directly or through colleagues.

\clause{3.3} \textbf{Still permitted}, all ordinary faculty duties: naming the
Company to the Technology Transfer Office as a potential licensee; technical
input on scope of work through the Research Administration Office's designated
contact; serving as named technical contact once an agreement is signed; and
signing the inventor's acknowledgment of a license or a royalty Distribution
Agreement.\footnote{Technology Transfer Office,
\href{https://techtransfer.byu.edu/resources}{``Working with the Technology
Transfer Office.''}}

\clause{3.4} \textbf{If the Company sells to BYU}, the sale happens only after
the Faculty Member is approved as an employee-vendor by Purchasing \& Travel,
Compensation, and Regulatory Accounting, and the Faculty Member takes no part
in the purchasing decision. An employee whose access to non-public BYU
information would give their company an unfair advantage may not do business
with the university at all.\footnote{Employee-Vendor Policy, in the sign-in
policy library at \href{https://policy.byu.edu}{policy.byu.edu}.}

\psection{Company-sponsored research and cost recovery}

\clause{4.1} Any Company use of lab personnel, students, or facilities runs
under a written agreement selected and signed by the authorized university
office: a \textbf{sponsored research agreement} for research work (the
default, at BYU's standard terms, with no lasting restrictions on
publication), a cost-recovery account for instrument time (4.4), or an
equipment agreement (6.4). No Company work happens in the lab outside a signed
scope of work. Only faculty on regular appointments and certain research
faculty may be principal investigators, so that responsibility cannot be
delegated to students, postdoctoral fellows, or staff.\footnote{\href{https://policy.byu.edu/view/intellectual-property-policy}{Intellectual
Property Policy}: such work ``must be administered as sponsored research
projects when there is a substantial use of university resources or when BYU
students are employed.''
\href{https://finserve.byu.edu/00000193-6566-d995-a9db-eff6a47a0001/sponsored-programs-handbook}{Sponsored
Projects Handbook} for the principal-investigator rule.}

\clause{4.2} \textbf{Independent review of the science.} A senior faculty
member with no Company interest, designated by the dean (\ph), reviews each
Company-sponsored project before work begins and at least annually after:
funding, student roles, publication freedom, and the separation of Company
from university work. The reviewer reports to the chair and the AAVP, and the
chair may pause affected work while a concern is open.

\clause{4.3} Publications and presentations from Company-sponsored or
license-related work disclose the Faculty Member's Company interest, and BYU's
stock position where BYU holds one. If BYU accepts Company stock, the
university's own interest is reviewed by the AAVP with the Office of General
Counsel, separately from this plan.

\clause{4.4} \textbf{Cost recovery.} Sales of instrument time outside a
research agreement run only through a cost-recovery account approved by the
chair, dean, and Regulatory Accounting, with outside users billed at the
external rate including overhead and the tax treatment settled before billing
starts.\footnote{\href{https://spo.byu.edu/00000190-3687-d777-a591-37af392d0001/rechargecenterpolicyprocedures-pdf}{Recharge
Center Policy and Procedures}. The published copy is dated 2005 and sits
outside the university policy library; nothing in this clause is settled until
Regulatory Accounting confirms which document governs.} That revenue is
budgeted to student wages, equipment, and student travel. None of it funds the
Faculty Member's salary, summer salary, or travel.

\psection{Students}

\clause{5.1} \textbf{Notice and an independent advocate.} Every lab member
receives written notice of the Faculty Member's Company interest, and BYU's if
BYU holds stock, on joining and annually after, with a copy to the chair. The
graduate coordinator or associate chair, who has no Company interest, meets
the students at least once a semester without the Faculty Member present, and
students may go to them, the chair, or the dean's office at any time.

\clause{5.2} \textbf{Participation is voluntary, and academic timelines
govern.} A comparable non-Company project, funding source, and supervisor are
available without penalty; retaliation is prohibited; and the choice never
affects grades, authorship, funding, or letters. Company tasks exist only
inside a signed scope of work, never as part of a BYU role. Thesis topics,
timelines, and publication schedules are set on academic grounds; no thesis is
embargoed and no defense placed under a confidentiality agreement. Where a
patent filing has to come first, the lab files a provisional application before
the defense.

\clause{5.3} \textbf{Hours and double employment.} University employment
respects BYU's caps: 20 hours a week in fall and winter semesters, with
stricter limits for international students, whose visa status can be at
stake.\footnote{Student Hiring and Employment Policy, and
\href{https://policy.byu.edu/view/conflict-of-time-or-interest-policy}{Conflict
of Time or Interest Policy} for the student disclosure trigger.} A student
working for BYU and the Company at once on related subject matter must
disclose it, and the advocate reviews the student's combined load each
semester.

\clause{5.4} \textbf{Company jobs and internships.} A student who wants to
work for the Company is referred to the advocate and the dean's office for
advice independent of the Faculty Member, and the advocate reviews the offer's
terms (duties, hours, pay, invention-assignment and confidentiality
obligations, visa implications) before it is
accepted.\footnote{\href{https://web.stanford.edu/group/OTL/documents/bp_faculty_sus.pdf}{Stanford's
Best Practices for Faculty Start-Ups} recommends the same step. BYU's
internship rules bar credit under a family member's supervision but are silent
on a faculty advisor's own company.} The Faculty Member does not supervise the
same student's thesis and Company work; a Company internship taken for credit
is arranged and graded by someone else. Any offer to a currently supervised
student requires an amendment first (section 10).

\psection{Company people and university resources}

\clause{6.1} \textbf{No university resources for Company work.} ``University
Employees may not use university resources, facilities, or work time to engage
in commercial activities, businesses, or `start-up'
companies.''\footnote{\href{https://policy.byu.edu/view/intellectual-property-policy}{Intellectual
Property Policy}. Using BYU networks for outside commercial purposes requires
written authorization from the responsible vice president and the chief
information officer.} All Company work runs on Company-owned laptops,
accounts, software, and storage, and on campus Company devices use their own
cellular connection rather than BYU networks. No Company data on BYU systems,
and no non-public BYU information on Company systems.

\clause{6.2} \textbf{Campus office.} Company use of the campus office is
limited to occasional calls and correspondence on Company devices: no Company
work performed there, and no Company records, operations, meetings, or
visitors. Written approval is requested from the \textbf{dean and the AAVP}
with the authorization of the outside role, and until it is granted there is
no Company use of the office at all.\footnote{\href{https://policy.byu.edu/view/conflict-of-time-or-interest-policy}{Conflict
of Time or Interest Policy} permits ``Incidental Use'' that is ``occasional or
infrequent'' and ``results in little or no cost to the university.'' The
Expectations of a Faculty Appointment Policy requires ``the explicit written
approval of the appropriate university supervisors'' to ``use his or her
campus office for business purposes.'' The 6.1 prohibition carries no chair-level waiver,
so a chair's signature alone does not clear it.}

\clause{6.3} \textbf{Company personnel.} Everyone working in the lab signs
BYU's standard participation and invention-assignment agreement. Company
employees never work in the lab informally or as
``volunteers.''\footnote{Volunteer Policy: service must be ``in furtherance of
the university's nonprofit objectives,'' ``unpaid service in furtherance of
the university's commercial activities'' is prohibited to comply with federal
law, and volunteers ``cannot, even on a temporary basis, displace employed
workers.''} They come on campus only under a written agreement, whether a visiting
appointment, the sponsored agreement, or a services contract, signed at the
tier its value and term require. Both sides acknowledge beforehand that work
done in the lab by a non-employee under faculty supervision can give BYU a
perpetual royalty-free license in the
result.\footnote{\href{https://policy.byu.edu/view/intellectual-property-policy}{Intellectual
Property Policy}: where such a person works ``under the supervision and
direction of a faculty member'' or ``uses substantial university resources,''
BYU may claim ``a non-exclusive, perpetual, royalty-free, paid-up, irrevocable
license to exploit, use, and sublicense the resulting Intellectual
Property.''}

\clause{6.4} \textbf{Equipment placements.} Company equipment enters the lab
only under a university-signed agreement, an equipment loan or a term of the
research agreement, settling ownership, insurance, maintenance, access,
export-control screening, and removal before delivery. The Faculty Member
signs nothing. Students use placed equipment for university research or within
the sponsored agreement, and it is reported as outside support on federal
applications.

\psection{Inventions and open-source release}

\clause{7.1} \textbf{Disclose first, release second.} The lab releases its
outputs as open-source software and hardware by default, but new inventions go
to the Technology Transfer Office \emph{before} anything is made public,
including a code repository, preprint, conference talk, or
thesis.\footnote{Public release first can permanently destroy foreign patent
rights.
\href{https://policy.byu.edu/view/intellectual-property-policy}{Intellectual
Property Policy}; Technology Transfer Office,
\href{https://techtransfer.byu.edu/resources}{``Preparing a Patent
Application.''}} Release proceeds once the office files a provisional application or releases
the rights. Review is bounded: normally 30 days, with up to 60 more only when
a patent filing is under way, and it never delays a defense, deposit, or
graduation.

\clause{7.2} \textbf{Open-source boundaries are set in advance}, in the lab's
published open-science policy and the Company's published licensing policy,
both adopted outside any BYU negotiation. BYU's ownership and approval rights
still govern every release.

\clause{7.3} \textbf{A clean line between BYU's inventions and the
Company's.} Both keep records of where, when, and with what resources each
invention was made. Because BYU's claim reaches inventions ``related to the current or
demonstrably anticipated business, research, or development of the
university,'' records alone do not settle ownership: a Company-side invention
near the lab's field is disclosed and the university's ownership decision
obtained in writing.\footnote{\href{https://policy.byu.edu/view/intellectual-property-policy}{Intellectual
Property Policy}. This plan cannot promise that a release or carve-out will be
granted.}

\clause{7.4} \textbf{Company assignment documents.} No Company
invention-assignment or prior-inventions agreement is signed by the Faculty
Member or any lab member without an express carve-out for BYU-owned and
BYU-assignable work. That form goes to every employee and intern, not only to
founders.

% ------------------------------------------------------------------- band C
\planband{C. How these conditions safeguard objectivity}{}

\psection{Why these conditions are enough}

{\footnotesize
\setlength{\extrarowheight}{2pt}
\begin{tabularx}{\textwidth}{@{} >{\RaggedRight\arraybackslash}p{1.6in} >{\RaggedRight\arraybackslash}X @{}}
\toprule
{\color{byunavy}\scshape What could bias a decision} &
{\color{byunavy}\scshape What removes it} \\
\midrule
Influence over a BYU and Company deal &
Total recusal on both sides (3.1, 3.2). Above ten percent, policy treats
influencing such negotiations as unmanageable, so the recusal is not a
judgment call. \\
\addlinespace
Payment tied to Company outcomes &
No salary, fees, bonuses, or honoraria (1.3); a forfeited royalty share
(1.5); travel reimbursed at cost, counted against the four-day allowance and
reported trip by trip (1.4). What remains is unsellable private stock, which no
single favorable term turns into money. \\
\addlinespace
Research shaped, suppressed, or delayed &
Independent review before work begins and annually after, and the chair may
pause the work (4.2). No embargo and no confidentiality on a defense (5.2),
bounded invention review that never delays graduation (7.1), and disclosure on
every publication (4.3). \\
\addlinespace
Pressure on students &
Real alternatives and no retaliation (5.2), an advocate outside the lab who
meets students alone (5.1), and independent review of any Company offer
(5.4). \\
\addlinespace
The two roles merging &
Company work on Company systems (6.1), personnel and equipment in the lab only
under signed agreements (6.3, 6.4), Company promises signed by a Company
officer (9.3). \\
\bottomrule
\end{tabularx}}

% ------------------------------------------------------------------- band D
\planband{D. How compliance is monitored}{}

\psection{Reporting, review, and remedies}

{\footnotesize
\setlength{\extrarowheight}{2pt}
\begin{tabularx}{\textwidth}{@{} >{\RaggedRight\arraybackslash}X >{\RaggedRight\arraybackslash}p{1.2in} >{\RaggedRight\arraybackslash}p{1.5in} @{}}
\toprule
{\color{byunavy}\scshape What} &
{\color{byunavy}\scshape When} &
{\color{byunavy}\scshape To whom} \\
\midrule
Outside-activity disclosure in
\href{https://facultyprofile.byu.edu/}{Faculty Profile}, time log (2.1),
advocate's student review (5.1), and independent reviewer's report (4.2) &
annually & chair, and the AAVP for the reviewer's report \\
Financial disclosure, including Company travel (1.4) & annually, updated
within 30 days & Research Administration Office \\
Company financing, ownership change, change of role or title, new proposed
transaction with BYU, or contemplated leave & within 30 days & chair and
Research Administration Office \\
\bottomrule
\end{tabularx}}

\vspace{2pt}
\clause{9.1} The chair monitors compliance through completion of the affected
projects, not only at annual
review.\footnote{\href{https://policy.byu.edu/view/financial-conflict-of-interest-in-research-policy}{Financial
Conflict of Interest in Research Policy}. The 30-day update of
\emph{financial} disclosures is required by that policy; the other 30-day
triggers above are terms of this plan.}

\clause{9.2} The AAVP reviews this plan annually with the chair and dean and
may strengthen it at any time using the tools policy names (an independent
monitor, removal from a project, reduction or elimination of the interest), or
end it in writing when the interests in section 1 no longer exist.

\clause{9.3} Commitments here about Company conduct are confirmed in a
one-page acknowledgment signed by an authorized Company officer and attached:
no use of BYU's name (1.7), no informal approaches to students (5.4, 6.3), and
no Company personnel or equipment in the lab outside signed agreements (6.3,
6.4). Failure to follow this plan is handled under the governing policies,
which provide for retrospective review.

% ------------------------------------------------------------------- band E
\planband{E. The investigator's agreement}{}

\psection{Changes that require amending this plan first}

Any officer title, employment relationship, or voting board role at the
Company; any leave of absence; any paid consulting for the Company, which
collides with taking part in research it sponsors and needs chair, dean, and
written AAVP approval (1.3); Company employment of a currently supervised
student; the Company selling anything to BYU (3.4); any Company agreement
asking the Faculty Member or a lab member to assign inventions or prior
inventions (7.4); any research involving human subjects that touches the
Company; and any change in the facts of section 1.

\psection{Approvals}

{\setlength{\extrarowheight}{2pt}\small
\begin{tabularx}{\textwidth}{@{} >{\RaggedRight\arraybackslash}p{2.25in} >{\RaggedRight\arraybackslash}X >{\centering\arraybackslash}p{1.5in} @{}}
\toprule
{\color{byunavy}\scshape Approver} &
{\color{byunavy}\scshape What the signature covers} &
{\color{byunavy}\scshape Signature / date} \\
\midrule
{[Sterling G.\ Baird]}, Faculty Member & agrees to be bound by this plan &
\sigline \\[0.3em]
\ph, Department Chair, Mechanical Engineering & authorizes the outside role;
monitors compliance (9.1) & \sigline \\[0.3em]
\ph, Dean, Ira A.\ Fulton College of Engineering & approves the
outside-activity conflict plan & \sigline \\[0.3em]
\ph, Associate Academic Vice President for Research \& Graduate Studies &
decides whether a conflict exists; approves the research conflict plan &
\sigline \\[0.3em]
\bottomrule
\end{tabularx}}

\vspace{1pt}
{\footnotesize\textbf{Routed for concurrence, not approval:} the Research
Administration Office, the Technology Transfer Office, Regulatory Accounting,
and the Office of General Counsel. Their approvals attach to the transactions
this plan governs, not to the plan itself.}

\end{document}
